\documentclass[11pt,letterpaper]{article}
\usepackage[T1]{fontenc}
\usepackage{newtxtext,newtxmath}
\usepackage[margin=0.8in,top=0.68in,bottom=0.7in]{geometry}
\usepackage{microtype}
\usepackage{xcolor}
\usepackage{enumitem}
\usepackage{fancyhdr}
\usepackage[hidelinks]{hyperref}
\definecolor{ink}{HTML}{18334A}
\setlength{\parindent}{0pt}
\setlength{\parskip}{5pt}
\setlength{\emergencystretch}{2em}
\pagestyle{fancy}
\fancyhf{}
\fancyfoot[L]{\small Working-style demo\enspace|\enspace Codex + Claude Code}
\fancyfoot[R]{\small\thepage\,/\,2}
\renewcommand{\headrulewidth}{0pt}
\renewcommand{\footrulewidth}{0.3pt}
\newcommand{\stage}[1]{\par\vspace{5pt}{\color{ink}\bfseries #1}\par\vspace{1pt}}
\newcommand{\prompt}[1]{\par\begingroup\leftskip=10pt\rightskip=6pt\small\textbf{My prompt:} ``#1''\par\endgroup}
\newcommand{\agent}[1]{\textbf{#1 --- work summary.} }
\hypersetup{pdftitle={Working-style demo: Codex and Claude Code},pdfsubject={Edited session excerpts from the string landscape research project}}
\begin{document}
{\Large\bfseries\color{ink} A working-style demo: Codex + Claude Code}\par
{\normalsize Building and auditing a string-compactification search}\par
{\small\itshape Edited from actual sessions. Prompt excerpts are verbatim; agent work is summarized.}
\vspace{3pt}

I used Codex and Claude Code to build a research codebase for heterotic string
compactifications with Standard-Model-like particle spectra. I set the scientific
scope, validation requirements and research priorities; the agents implemented
code, ran experiments and developed mathematical checks. The project began on a
Mac mini M4 with 16\,GB of memory.

\stage{1. Give Codex a specification with scientific acceptance criteria}
\prompt{The first goal is correctness and reproducibility. Performance comes second.}

My initial specification restricted the problem to five-line-bundle sums on
favourable Calabi--Yau threefolds. It required primary-source checks, exact
arithmetic, data provenance, a slow reference enumerator and a separate optimized
search. Every filter had to state whether it was exact, necessary, sufficient or
heuristic. An unsuccessful numerical search could not become a proof of absence.

\agent{Codex} Built the topology and bundle-algebra core, tests, CLI,
checkpoints and certificate verification. It reproduced a published $X_{7447}$
model's two indices, both $-12$, and matching second Chern classes. A sign
inconsistency in a published formula was checked against another primary source,
the Chern-class derivation and the known model before being used in a filter.
The initial baseline passed 71 tests.

\stage{2. Hand the existing repository to Claude Code for a stronger audit}
\prompt{The reference implementation must not reuse the optimized hierarchical pruning logic.}
\prompt{Require exact set equality.}

Once the first benchmark matched the published count of five models, I asked
Claude Code to inspect the existing repository and continue from its validated
state. I required an independent replay of all 43,046,721 ordered assignments
in the smallest tetraquadric box. I also requested the actual published catalog,
so the comparison would establish model-for-model agreement after the allowed
permutations, rather than merely matching totals.

\agent{Claude Code} Implemented the independent replay and confirmed the
same empty survivor set. It downloaded and parsed the catalog, then established
exact orbit-set equality at the next benchmark. Switching tools was not itself
an independent check; the separately implemented enumeration and explicit set
comparison supplied that evidence.

\stage{3. Ask what a pruning rule actually proves}
\prompt{A branch may be pruned only when the target is mathematically impossible for every completion inside the declared coefficient box.}

My follow-up required a mathematical statement, proof or conjectural label,
counterexample search, differential test and performance measurement. This led
to an audit of the finite-coefficient slope test. Its rejection witnesses were
sound, but passing the test did not establish stability. At the next charge
bound, it retained an extra orbit. Claude Code implemented a complete version
of that obstruction test, rejecting the extra orbit and restoring exact
agreement with the seven published orbits in that box.

\newpage
{\large\bfseries\color{ink} Steering the work beyond a successful benchmark}\par

\stage{4. Change priorities when more scanning stops answering the main question}
\prompt{have you done phase 8 and 9?}

After the search and catalog comparisons worked, I redirected Claude Code
toward the physical spectrum on quotient geometries. I asked it to compute the
symmetry action on cohomology rather than rely on database lookups, including
the vector-like sector needed to count Higgs pairs. The resulting work extended
the code to equivariant cohomology and Wilson-line projections. This changed the
next task from enlarging a charge box to testing what the surviving models
actually predict.

\stage{5. Audit additional candidates before claiming a discovery}
\prompt{Do not call any candidate new until this audit is complete.}
\prompt{Do not infer a coupling merely from U(1) charge conservation.}

When eleven additional orbits passed the implemented MSSM-spectrum conditions,
I requested an immutable candidate manifest, comparison against the original
scan's full conditions, and an explanation of each absence from its database.
I then asked for geometric Yukawa calculations and dangerous-operator checks,
with published models processed through the same pipeline as controls.

\agent{Claude Code} Implemented the audit and supported cohomological
coupling calculations. The recorded ranking classified nine additional orbits
as phenomenologically problematic and two as unknown under the stated
perturbative criteria. None met the pipeline's strong-candidate requirements.
The historical omissions were explained by restrictions of the earlier scan.
A larger list of candidates was therefore not presented as eleven established
new Standard Models.

\stage{6. Keep corrections visible and strengthen the tests}

The sessions also contain corrections. During the Claude Code work, a
cohomology comparison initially checked only single-degree cases and gave an
overly broad agreement report. Checking every degree exposed apparent
character-label disagreements. Further investigation traced these to applying
an equivariant-structure shift on one side of the comparison but not the other.
The comparison was corrected, the regression requirement strengthened, and the
explanation appended to \texttt{RESULTS.md}.

A separate correction changed unsupported higher-level coupling calculations
from ``zero'' to \texttt{UNKNOWN}. These distinctions matter: an uncomputed
coupling and a mathematically vanishing coupling lead to different physical
conclusions. My specification required the notebook to retain failures and
unresolved cases, not just successful benchmarks.

\stage{What this demonstrates about how I work}

I start with a precise research question and define what evidence would count
as progress. I give agents room to implement and experiment, then use
independent replay, catalog comparison and explicit mathematical assumptions to
challenge their outputs. I redirect the work when the bottleneck changes, and
keep promising candidates separate from verified claims. The accompanying
repository contains the code, tests, benchmark artifacts, mathematical notes
and research log needed to inspect this process.

\vfill
{\footnotesize\color{ink}\textbf{Source trail.} Initial Codex build session and master prompt;
local Claude Code session \emph{String landscape exhaustive replay and verification};
repository \texttt{RESULTS.md}, slope-obstruction notes and candidate-ranking artifacts.
}
\end{document}
